% 第2章 线性直流电路（章节子文件，需被主文件调用）
\chapter{线性直流电路}

\section{电阻网络的等效}
电阻的串联

电阻的并联

电阻的星形联结和三角形联结

电阻的星形联结，简称$\mathrm{Y}$形联结或$\mathrm{T}$形联结，。三角形联结，。

可归纳为
\[
R_{1,2,3}=\frac{\boldsymbol{\Delta}\text{形联结中与该节点相连两电阻之积}}{\boldsymbol{\Delta}\text{形联结中各电阻之和}}
\]

可归纳为
\[
R_{1,2,3}=\frac{\mathrm{Y}\text{形联结各电阻两两乘积之和}}{\mathrm{Y}\text{形联结中不相邻电阻}}
\]

\section{含源支路的等效变换}
图表示，它们的电压、电流关系分别为
\[
U=U_s-R_iI
\]
\[
I=I_s-G_iU
\]

当电路参数满足一定条件时，电压源和电阻串联支路与电流源和电导并联支路可以相互等效。为求出等效条件，将式改写成
\[
U=\frac{1}{G_i}I_s-\frac{1}{G_i}I
\]
比较式与式得到两电路的等效变换条件为
\[
U_s=\frac{I_s}{G_i},\quad R_i=\frac{1}{G_i}
\]
图中电压源

受控电压源与电阻的串联也可等效成受控电流源与同一电阻的并联，或反之

\section{支路电流法}

\section{回路电流法}
回路电流

回路电流方程的列写

\section{节点电压法}

\section{运算放大器}

\section{含运算放大器电路的分析}
反相放大器

同相放大器

加法器

差分放大器